El desarrollo del proyecto permitió obtener un Analizador Vectorial de Redes de un puerto para la banda UHF de 1,8 a 2,2 GHz funcional, de bajo costo y tamaño compacto, orientado a la medición no invasiva de glucemia mediante sensores de microondas. Se diseñó, construyó, calibró y ensayó a nivel de hardware y software, \textcolor{magenta}{y permite realizar mediciones comparables con las de un dispositivo comercial}. El prototipo desarrollado integra la generación, reflectometría y heterodinación de señales en un módulo portable que satisface los requerimientos fijados para la aplicación. Por lo tanto, se puede afirmar que el objetivo general fue logrado. 

A partir de los ensayos realizados y el análisis de los resultados obtenidos, y en base a los objetivos particulares, se concluye que el sistema:
\begin{itemize}
    \item Fue implementado en una placa de circuito impreso de 4 capas, lo que permitió tener un plano de masa continuo bajo las líneas de transmisión coplanares y optimizar la distribución de la alimentación, reduciendo posibles fuentes de error y ruido electromagnético y garantizando la adaptación de impedancias;
    \item Fue capaz de generar un barrido en la banda de frecuencias de interés \textcolor{magenta}{y esa banda pudo ser extendida, lo que permite extender las posibilidades a futuro de diseñar sensores que resuenen en otras frecuencias};
    \item Implementó un sistema heterodino basado en dos sintetizadores de RF (ADF4351) y un mezclador de dos canales (ADL5802) que demostró ser eficaz para trasladar la información desde la banda de microondas a una frecuencia intermedia, que simplifica la etapa de acondicionamiento de las señales y la digitalización, reduciendo el costo del hardware de adquisición;
    \item Incluyó un acoplador direccional basado en un puente resistivo con una directividad superior a 36,5 dB en toda la banda de trabajo, superando ampliamente la directividad mínima de 15 dB planteada en la Especificación de Requerimientos \ref{apend:ER}, asegurando una adecuada separación entre las ondas incidente y reflejada a costa de una pérdida por inserción mayor que la de un circuito comercial;
    \item Implementó una rutina de calibración SOL de tres términos de error, que eliminó los errores sistemáticos garantizando la repetibilidad en la determinación del coeficiente de reflexión (S11).
\end{itemize}

Todos estos objetivos permitieron obtener exitosamente el coeficiente de reflexión del sistema bajo prueba, habiéndose medido cargas como cortocircuitos, circuitos abiertos, resistencias, cables y sensores de reflexión. En cuanto a estos últimos, se detectaron desplazamientos en la frecuencia de resonancia asociados a las variaciones de los parámetros eléctricos del medio, demostrando que el $\epsilon-\mu VNA$ es una herramienta útil para la línea de investigación en glucometría no invasiva del Laboratorio de Comunicaciones. 

\section{Trabajos futuros}
A partir de los resultados obtenidos, se proponen como líneas de trabajo a futuro:

\begin{itemize}
    \item Diseñar un gabinete que permita contener el VNA y transportarlo fácilmente;
    \item Embeber el software y firmware en un microcontrolador que permita independizar el dispositivo de una PC y prescindir del uso de cables puente, sensibles a radiación electromagnética; 
    \item En línea con los puntos anteriores, agregar una forma de visualizar los resultados ya sea con una pantalla incorporada o en una aplicación móvil mediante Bluetooth o Wi-Fi;
    \item Optimizar el tiempo de barrido calculando el tiempo de estabilización del enganche y adquiriendo la menor cantidad de tiempo posible;
    \item Explorar la posibilidad de incorporar un segundo canal para medir transmisión (parámetro $S_{21}$) que permita caracterizar otro tipo de sensores;
    \item Seguir miniaturizando el PCB e incorporar blindaje electromagnético para minimizar el ruido.
\end{itemize}